\section{Step 64 Build the alphabet that describes the data alphabets}\label{png-huff-s64}
Fill \texttt{huff\_code\_tree}. Call the shared builder with all nineteen code-length slots, its table at frame offset 928 and strict policy zero. Then establish the combined data-length cursor at frame offset 1096 and its end after literal/length count plus distance count. Retain cursor and end at offsets 24 and 16.

This fixture has 259 combined entries. Its code-length table contains symbols 0, 1, 2, 16, 17 and 18, ordered by lengths one, two, three, four, five and five. Their decoded values now describe code lengths or repeats; they are not inflated row bytes.

\paragraph{Observe} Use \texttt{s64-1}; stop before \texttt{huff\_code\_tree\_ready}. Inspect the state described above before continuing.

\Needspace{7\baselineskip}
\paragraph{Check} What does this intermediate Huffman table decode?
\begin{itemize}
\item[A.] Packed display words.
\item[B.] Only backward distances.
\item[C.] Lengths and repeat operators for the two data alphabets.
\end{itemize}
